% Method — full prose v1.

\subsection{Setup}

We study two small open-weight models with distinct architectures and
training pipelines --- Gemma-4 E4B (\texttt{google/gemma-4-E4B-it}) and
Ministral-3 3B (\texttt{mistralai/Ministral-3-3B-Reasoning-2512}, an
explicitly reasoning-tuned release) --- run locally in thinking mode, in
4-bit MLX quantization (\texttt{mlx-community} conversions), with full
token-level instrumentation (per-token logits and final-layer hidden
states). Problems come from the MATH benchmark \citep{hendrycks2021math}.
Because breakthrough measurement is only informative on problems a model
sometimes solves, cohorts were selected by outcome-blind
intermediate-difficulty screens, frozen before any probe outcome existed: a
development cohort (20 problems), a supplement cohort (20), and an expansion
cohort (49) whose screen was repaired after a recorded gate failure
(App.~\ref{app:amendments}). Together they yield 89 problems and $89 \times
2 = 178$ problem--model cells. Each cell has one instrumented base
trajectory generated with a 16{,}384-token budget. Research splits (54
train / 18 validation / 17 test problems) were assigned before any outcome
was observed; the test split is spent exactly once, in
Section~\ref{sec:results-prediction}.

\subsection{Truncation probe and the budget-fit time \Tf}

The probe follows the truncation-and-resampling tradition
\citep{lanham2023,bogdan2025anchors}: truncate the model's \emph{own}
trajectory at an anchor $t$ (log-spaced from 16 to 8{,}192 tokens), and
sample $m=4$ continuations from the truncated prefix under a fixed budget
of $B = 1{,}024$ reasoning tokens plus a 512-token answer reserve, with
deterministic per-branch seeds. The per-anchor solve rate
$\phat(t; B)$ estimates the value of the partial state. A
\emph{breakthrough} is the first anchor with $\phat \ge \tauthr$
($\tauthr = 0.75$) that remains above threshold at the next anchor;
crossings are refined by bisection and recorded as intervals, and
trajectories whose curve never crosses are right-censored at their final
anchor. We call the resulting time \Tf$(B)$, the \emph{budget-fit time}:
the point at which a solution first fits the continuation budget. As
Section~\ref{sec:results-regimes} shows, \Tf{} is what an uncontrolled
probe reports as a breakthrough --- and it conflates two very different
things.

\subsection{Restart control and the prefix-value time \Tv}

To separate ``the prefix carries value'' from ``the budget became
sufficient,'' we measure each problem's \emph{restart curve} $R(C)$: the
solve rate of from-scratch attempts (empty prefix, same prompt) at budgets
$C \in \{1{,}024, 2{,}048, 4{,}096, 8{,}192\}$, four attempts each, with
$\Rhat$ interpolated log-linearly between grid points. The
\emph{advantage} of a prefix at matched total compute is
$\mathrm{adv}(t) = \phat(t; B) - \Rhat(t + B)$, and the prefix-value time
$\Tv(\delta)$ is the earliest stable anchor with $\phat \ge \tauthr$ and
$\mathrm{adv} \ge \delta$ (primary margin $\delta = 0.5$; $0.25$ and
$0.75$ reported as sensitivity, along with a conservative variant that
takes the upper envelope of the bracketing grid values of $\Rhat$). Cells
are then classified by frozen precedence: \emph{instant} (event interval
upper bound $\le 16$ tokens), \emph{budget-limited} ($\Rhat(4{,}096) \ge
\tauthr$: a restart solves it), \emph{prefix-limited} (a $\delta = 0.5$
crossing exists), \emph{no-crossing} (\Tf{} exists but no advantage
crossing), \emph{terminal} (a replicated event at the final anchor with no
stability anchor; annotated, never promoted to \Tv), and \emph{unsolved}.

\subsection{Noise hardening}

With $m = 4$, a $3/4$ crossing occurs with probability $0.31$ under a true
success rate of $0.5$, so single-shot labels are optimistic. Two frozen
rules harden them. First, threshold-censored trajectories (final-anchor
rate $\ge \tauthr$, no stability anchor available) get four additional
branches, and a terminal event is recorded only if the pooled rate reaches
$\ge 6/8$ (under $p = 0.5$ this has probability $0.14$). Second, every
ambiguous cell ($1$--$3$ successes of $4$) across the expansion cohort was
enlarged to eight attempts before labels were derived (148 cells). Both
rules cut in both directions in practice: replication confirmed 13 of 16
candidates for one model and rejected the other model's single candidate,
and pooling removed six events while adding one.

\subsection{Pre-registration discipline}

Every decision rule above was frozen in written, committed amendments
before the outcomes it governs were observed; gate failures (two cohort
gates and one pilot gate) are recorded as failures and resolved by
amendment, never by relabeling. The confirmatory prediction protocol of
Section~\ref{sec:results-prediction} --- endpoints, feature sets, model
class, power gates, and success criterion --- was committed before the
single test-split evaluation, and every post-hoc analysis is labeled as
such in the amendment that reports it. Appendix~\ref{app:amendments} gives
the full timeline.
